Let @@M0@@ be a closed subscheme of @@M1@@, and @@M2@@ the open complement of @@M3@@ in @@M4@@. Denote by @@M5@@ the restriction of @@M6@@ over @@M7@@, and by @@M8@@ the restriction of @@M9@@ over @@M10@@. Finally, still denote by @@M11@@ the images of the @@M12@@ in @@M13@@ and @@M14@@, and denote by:
@@M15@@
the homomorphisms defined in the same way as @@M16@@.

Consider the diagram
\sgaStableTarget{sga6.e09.eq.1-2-1}@@M17@@
where @@M18@@ and @@M19@@ denote the immersions of @@M20@@ and @@M21@@ into @@M22@@, and @@M23@@ and @@M24@@ the immersions of @@M25@@ and @@M26@@ into @@M27@@. Let us show that if @@M28@@ and @@M29@@ are surjective, the same then holds for @@M30@@.

The two horizontal rows of the diagram are exact by 1.1. One then sees that to show that @@M31@@ is surjective, it suffices to show that the two squares commute.

Let @@M32@@; we have:
@@M33@@
@@M34@@
the last equality following from the projection formula IV 2.11.1.2. The equality @@M35@@ therefore follows from the equality @@M36@@, which holds by IV 3.1.1 since @@M37@@ is
proper and @@M38@@ flat.
